\documentclass[11pt,a4paper]{article}
\usepackage{curso}
\begin{document}
\portada{Ejercicios --- Lección}{2.8}{Procesos de saltos de Markov en tiempo continuo}

\section{Ejercicios}
\begin{enumerate}[label=\textbf{E\arabic*.}]
\item (B) Escriba $Q$ para el pbit y resuelva $p(t)$ desde $p(0)=(1,0)$.
\item (B) Muestre que Poisson$(k/\gamma)$ es estacionaria para producción $k$ y degradación $\gamma n$.
\item (B) Media y varianza del tiempo para completar $N=100$ saltos de tasa $k=10^6$ s$^{-1}$.
\item (B) Describa los dos sorteos de Gillespie.
\item (I) Caminata con saltos $+1$ (tasa $a$) y $-1$ (tasa $b$), $a>b$, desde 0 hasta $N$: muestre $\langle\tau\rangle=N/(a-b)$.
\item (I) Verifique balance detallado para el proceso de nacimiento y muerte con la Poisson.
\item (A) Implemente Gillespie en JAX para el pbit y estime $\tau_c$ de la autocorrelación.
\end{enumerate}

\section{Soluciones}
\begin{solucion}
\textbf{E1.} $Q=\begin{pmatrix}-k_+&k_+\\k_-&-k_-\end{pmatrix}$; $p_{+}(t)=\frac{k_+}{k_++k_-}\big(1-e^{-(k_++k_-)t}\big)$, $p_-=1-p_+$.

\textbf{E2.} Balance de flujos entre $n$ y $n+1$: $\pi_nk=\pi_{n+1}\gamma(n+1)\Rightarrow\pi_{n+1}/\pi_n=\frac{k/\gamma}{n+1}$, que es Poisson.

\textbf{E3.} $\langle\tau\rangle=10^{-4}$ s; $\Var=10^{-10}$ s$^2$ ($\sigma=10\ \mu$s, 10\,\% de la media).

\textbf{E4.} Tiempo $\sim\mathrm{Exp}(R)$; destino con probabilidad $Q_{ij}/R$.

\textbf{E5.} $m_n$ cumple $a(m_{n+1}-m_n)+b(m_{n-1}-m_n)=-1$; con deriva positiva, $m_n=(N-n)/(a-b)$ (despreciando la frontera inferior si es reflectante y lejana); desde $0$: $N/(a-b)$.

\textbf{E6.} Es la misma ecuación de E2: $\pi_nQ_{n,n+1}=\pi_{n+1}Q_{n+1,n}$.

\textbf{E7.} (Abierto.) Debe obtenerse $\tau_c\approx1/(k_++k_-)$.
\end{solucion}
\end{document}
